L'Estació Espacial Internacional (EEI) orbita a 400~km per sobre la superfície terrestre i té la seva fi planificada per a l'any 2031.

\begin{apartat}{1,25}
A partir de la llei de gravitació universal, deduïu l'expressió de la velocitat orbital en funció del radi orbital. Calculeu la velocitat de l'EEI en òrbita i el nombre de voltes a la Terra que fa cada dia.
\end{apartat}

\begin{apartat}{1,25}
De manera gradual i controlada s'abaixarà l'òrbita fins als 280~km d'altura per sobre la superfície terrestre. Calculeu l'energia mecànica de l'EEI en aquesta òrbita i justifiqueu-ne el signe. Posteriorment, l'estació caurà des d'aquesta altura al mig de l'oceà Pacífic. Calculeu amb quina energia cinètica impactaria l'estació contra l'aigua sense tenir en compte els efectes de l'atmosfera terrestre.
\end{apartat}

\dades{%
  $G = 6,67\cdot 10^{-11}\ \mathrm{N\,m^2\,kg^{-2}}$\\
  Massa de la Terra: $M_\mathrm{T} = 5,98\cdot 10^{24}\ \mathrm{kg}$\\
  Radi de la Terra: $R_\mathrm{T} = 6,37\cdot 10^{6}\ \mathrm{m}$\\
  Massa de l'EEI $= 430\cdot 10^{3}\ \mathrm{kg}$}
